% MLN122 · Chương 3 — Hàng hoá sức lao động
% Thiết kế: LLCTFPT v2.0 (khóa). Biên dịch bằng XeLaTeX (2 lần).
\documentclass[aspectratio=169,11pt]{beamer}
\usetheme[base=Warsaw]{LLCTFPT}

\PTTESession[MLN122 · Hàng hoá sức lao động]
  {Chương 3: Giá trị thặng dư trong nền kinh tế thị trường\\I. Lý luận của C.Mác về giá trị thặng dư -- 1. Nguồn gốc của giá trị thặng dư}
  {HÀNG HOÁ SỨC LAO ĐỘNG}
\author[ThS. Nguyễn Trung Hiếu]{ThS. Nguyễn Trung Hiếu}
\institute[ĐH FPT]{BỘ MÔN LÝ LUẬN CHÍNH TRỊ \& KỸ NĂNG MỀM}
\date[Fall 2026]{Fall 2026}

\begin{document}
%% PTTE-ID: S3-01
\PTTECoverFrame
%% PTTE-ID: S3-02
\PTTEContentsFrame

\section[Công thức chung của tư bản]{a. Công thức chung của tư bản và mâu thuẫn của nó}

%% PTTE-ID: S3-03
\begin{frame}{Từ công thức chung của tư bản đến câu hỏi cần giải đáp}{a. CÔNG THỨC CHUNG CỦA TƯ BẢN}
\begin{PTTEPoints}{CÔNG THỨC T -- H -- T'}
\PTTEPoint{Công thức}{Tiền được ứng ra để mua hàng hoá, rồi bán hàng hoá để thu về nhiều tiền hơn: T -- H -- T', trong đó T' = T + $\Delta$T.}
\PTTEPoint{Giá trị thặng dư}{Số tiền trội ra $\Delta$T so với số tiền ứng ra ban đầu được gọi là giá trị thặng dư (m).}
\PTTEPoint{Tư bản}{Giá trị mang lại giá trị thặng dư là tư bản; mục đích của lưu thông tư bản là giá trị lớn lên.}
\end{PTTEPoints}
\begin{CQ}Giá trị thặng dư từ đâu mà có?\end{CQ}
\end{frame}

%% PTTE-ID: S3-04
\begin{frame}{Mâu thuẫn của công thức chung của tư bản}{a. CÔNG THỨC CHUNG CỦA TƯ BẢN}
\begin{PTTEPoints}{HAI KHẢ NĂNG TRONG LƯU THÔNG}
\PTTEPoint{Trao đổi ngang giá}{Chỉ thay đổi hình thái giá trị (hàng thành tiền, tiền thành hàng); tổng giá trị không đổi nên không có giá trị thặng dư.}
\PTTEPoint{Trao đổi không ngang giá}{Người này được lợi thì người kia thiệt; tổng giá trị toàn xã hội vẫn không tăng.}
\PTTEPoint{Ngoài lưu thông}{Nếu người sản xuất chỉ đứng riêng một mình, giá trị chỉ được bảo tồn hoặc chuyển sang sản phẩm mới, không tự lớn lên.}
\end{PTTEPoints}
\begin{KeyClaim}Giá trị thặng dư vừa không thể xuất hiện ở ngoài lưu thông, vừa không thể chỉ do bản thân lưu thông sinh ra. Phải tìm một hàng hoá đặc biệt: hàng hoá sức lao động.\end{KeyClaim}
\end{frame}

\section[Hàng hoá sức lao động]{b. Hàng hoá sức lao động}

%% PTTE-ID: S3-05
\begin{frame}{Sức lao động và hai điều kiện để trở thành hàng hoá}{b. HÀNG HOÁ SỨC LAO ĐỘNG}
\begin{Concept}Sức lao động là toàn bộ những năng lực thể chất và tinh thần tồn tại trong một con người, được người đó sử dụng vào sản xuất hàng hoá.\end{Concept}
\begin{PTTEPoints}{HAI ĐIỀU KIỆN ĐỂ SỨC LAO ĐỘNG TRỞ THÀNH HÀNG HOÁ}
\PTTEPoint{Điều kiện 1}{Người lao động phải được tự do về thân thể, làm chủ được sức lao động của mình và có quyền bán sức lao động như một hàng hoá.}
\PTTEPoint{Điều kiện 2}{Người lao động không có đủ tư liệu sản xuất cần thiết để tự kết hợp với sức lao động của mình, nên buộc phải bán sức lao động để sống.}
\end{PTTEPoints}
\end{frame}

%% PTTE-ID: S3-06
\begin{frame}{Sức lao động trở thành hàng hoá: ý nghĩa lịch sử}{b. HÀNG HOÁ SỨC LAO ĐỘNG}
\begin{PTTEPoints}{ĐIỀU KIỆN CẦN VÀ ĐỦ}
\PTTEPoint{Thiếu điều kiện 1}{Người lao động không được tự do về thân thể thì không thể đem bán sức lao động của chính mình.}
\PTTEPoint{Thiếu điều kiện 2}{Người lao động có tư liệu sản xuất thì tự sản xuất, không cần bán sức lao động.}
\PTTEPoint{Đủ cả hai}{Sức lao động mới trở thành hàng hoá; đây là điều kiện để tiền biến thành tư bản.}
\end{PTTEPoints}
\begin{Takeaway}Sức lao động chỉ trở thành hàng hoá khi có đồng thời cả hai điều kiện.\end{Takeaway}
\end{frame}

%% PTTE-ID: S3-07
\begin{frame}{Hai thuộc tính của hàng hoá sức lao động}{b. HÀNG HOÁ SỨC LAO ĐỘNG}
\begin{PTTEPoints}{NHƯ MỌI HÀNG HOÁ, SỨC LAO ĐỘNG CÓ HAI THUỘC TÍNH}
\PTTEPoint{Giá trị}{Do thời gian lao động xã hội cần thiết để sản xuất và tái sản xuất sức lao động quyết định.}
\PTTEPoint{Giá trị sử dụng}{Thể hiện ở công dụng của nó, tức là quá trình người lao động làm việc cho người mua.}
\end{PTTEPoints}
\begin{KeyClaim}Hàng hoá sức lao động có tính chất đặc biệt: giá trị sử dụng của nó có khả năng tạo ra một lượng giá trị lớn hơn giá trị của bản thân nó.\end{KeyClaim}
\end{frame}

%% PTTE-ID: S3-08
\begin{frame}{Giá trị của hàng hoá sức lao động}{b. HÀNG HOÁ SỨC LAO ĐỘNG}
\begin{PTTEPoints}{GIÁ TRỊ ĐƯỢC ĐO GIÁN TIẾP QUA GIÁ TRỊ TƯ LIỆU SINH HOẠT, GỒM BA BỘ PHẬN}
\PTTEPoint{Nuôi sống bản thân}{Giá trị các tư liệu sinh hoạt cần thiết để tái sản xuất sức lao động của người lao động.}
\PTTEPoint{Nuôi sống gia đình}{Giá trị các tư liệu sinh hoạt cần thiết cho gia đình người lao động, bảo đảm sự kế tục của lực lượng lao động.}
\PTTEPoint{Đào tạo}{Phí tổn đào tạo người lao động để có được trình độ, kỹ năng nhất định.}
\end{PTTEPoints}
\begin{Takeaway}Giá trị sức lao động gồm cả yếu tố vật chất, yếu tố tinh thần và yếu tố lịch sử; nó thay đổi theo từng nước, từng thời kỳ.\end{Takeaway}
\end{frame}

%% PTTE-ID: S3-09
\begin{frame}{Giá trị sử dụng của hàng hoá sức lao động}{b. HÀNG HOÁ SỨC LAO ĐỘNG}
\begin{PTTEPoints}{ĐẶC ĐIỂM CỦA GIÁ TRỊ SỬ DỤNG}
\PTTEPoint{Thể hiện trong quá trình lao động}{Người mua chỉ nhận được giá trị sử dụng khi tiêu dùng sức lao động, tức là khi người lao động làm việc.}
\PTTEPoint{Tạo ra giá trị mới}{Trong quá trình lao động, sức lao động tạo ra một lượng giá trị mới lớn hơn giá trị của bản thân nó.}
\PTTEPoint{Phần dôi ra}{Phần giá trị lớn hơn đó chính là giá trị thặng dư mà nhà tư bản chiếm đoạt.}
\end{PTTEPoints}
\begin{Takeaway}Giá trị sử dụng của sức lao động là nguồn gốc của giá trị thặng dư; đó là chìa khoá giải đáp mâu thuẫn của công thức chung của tư bản.\end{Takeaway}
\end{frame}

%% PTTE-ID: S3-10
\begin{frame}{Phân biệt sức lao động và lao động}{b. HÀNG HOÁ SỨC LAO ĐỘNG}
\begin{PTTEPoints}{HAI KHÁI NIỆM KHÁC NHAU}
\PTTEPoint{Sức lao động}{Là năng lực lao động, tồn tại trong con người; là cái được đem ra mua bán, nên là hàng hoá.}
\PTTEPoint{Lao động}{Là sự tiêu dùng sức lao động trong thực tế, là hoạt động tạo ra giá trị; bản thân lao động không phải là hàng hoá.}
\end{PTTEPoints}
\begin{Misconception}Nhà tư bản mua lao động của công nhân. Thực tế, nhà tư bản mua sức lao động; lao động là quá trình sử dụng sức lao động đã mua.\end{Misconception}
\end{frame}

%% PTTE-ID: S3-11
\begin{frame}{Tổng kết và câu hỏi thảo luận}{b. HÀNG HOÁ SỨC LAO ĐỘNG}
\begin{Takeaway}Hàng hoá sức lao động là hàng hoá đặc biệt: giá trị sử dụng của nó là nguồn gốc của giá trị thặng dư. Nhờ đó, tiền biến thành tư bản mà không vi phạm quy luật giá trị.\end{Takeaway}
\begin{CQ}Vì sao nói người lao động làm thuê vừa “tự do”, vừa buộc phải bán sức lao động? Giá trị sử dụng của sức lao động khác gì so với các hàng hoá thông thường?\end{CQ}
\end{frame}

\end{document}
